\section{Generate Parentheses}
\label{sec:rec-generate-parentheses}

\problemstatement{Given an integer $n$, generate all combinations of
\textbf{well-formed} (valid) parenthesis strings using exactly $n$ pairs
of parentheses.}

\begin{patternbox}
\emph{``Generate all valid \ldots''} signals backtracking with
\textbf{validity-aware pruning}: rather than generating all $\binom{2n}{n}$
arrangements of $n$ open and $n$ close parentheses and filtering out the
invalid ones afterward (as in the Stack and Queue chapter's Valid
Parentheses problem \emph{checking} a given string), we instead only ever make a choice that
keeps the string-so-far on track to \emph{possibly} become valid,
pruning illegal branches the moment they would become unsalvageable.
\end{patternbox}

\subsection*{Understanding the problem carefully}

At any point while building the string, two counts matter: how many
open parentheses have been placed so far, and how many close parentheses
have been placed so far. An open parenthesis can be placed as long as
fewer than $n$ have been used so far (we must not exceed $n$ total). A
close parenthesis can be placed only if doing so would not create more
closes than opens at any prefix -- i.e.\ only if the current count of
closes is still strictly less than the current count of opens (otherwise
an unmatched close would appear, which the Stack and Queue chapter's Valid Parentheses problem
already established invalidates a parenthesis string).

\begin{ideabox}[Two Guarded Choices Instead of a Binary Pick/Not-Pick]
$\textsc{Solve}(\textit{current}, \textit{openCount}, \textit{closeCount})$:
if $|\textit{current}| = 2n$: record $\textit{current}$; return.
Otherwise: if $\textit{openCount} < n$: append \texttt{'('}; recurse with
$\textit{openCount}+1$; remove (undo). If $\textit{closeCount} <
\textit{openCount}$: append \texttt{')'}; recurse with
$\textit{closeCount}+1$; remove (undo).
\end{ideabox}

\subsection*{Why these two guards are exactly the validity conditions, applied early}

This is the backtracking analogue of
the Stack and Queue chapter's Valid Parentheses stack-based validity check,
applied \emph{during construction} rather than after the fact: a
parenthesis string is valid exactly when every prefix has at least as
many opens as closes, and the total opens and closes both equal $n$. The
guard $\textit{closeCount} < \textit{openCount}$ is precisely
``this prefix currently has more opens than closes, so one more close is
still safe''; the guard $\textit{openCount} < n$ simply caps the total
opens at $n$. Because these guards are checked \emph{before} a character
is ever placed, no invalid prefix is ever constructed in the first
place -- the search space explored is exactly the valid strings (and
their valid-so-far prefixes), never the much larger space of all
arrangements.

\subsection*{Algorithm}

\begin{enumerate}
  \item $\textsc{Solve}(\textit{current}, \textit{openCount}, \textit{closeCount})$:
  \begin{enumerate}
    \item If $|\textit{current}| = 2n$: record $\textit{current}$;
          return.
    \item If $\textit{openCount} < n$: append \texttt{'('}; recurse;
          remove.
    \item If $\textit{closeCount} < \textit{openCount}$: append
          \texttt{')'}; recurse; remove.
  \end{enumerate}
  \item Call $\textsc{Solve}(\texttt{""}, 0, 0)$.
\end{enumerate}

\subsection*{Dry run}

\dryrun{Take $n=2$.}

\begin{itemize}
  \item $\textsc{Solve}(\texttt{""},0,0)$: open allowed ($0<2$):
        append \texttt{(}. $\textsc{Solve}(\texttt{"("},1,0)$.
  \begin{itemize}
    \item Open allowed ($1<2$): append \texttt{(}.
          $\textsc{Solve}(\texttt{"(("},2,0)$.
    \begin{itemize}
      \item Open not allowed ($2\not<2$). Close allowed ($0<2$): append
            \texttt{)}. $\textsc{Solve}(\texttt{"(()"},2,1)$: close
            allowed ($1<2$): append \texttt{)}.
            $\textsc{Solve}(\texttt{"(())"},2,2)$: length $4=2n$,
            record \texttt{"(())"}.
    \end{itemize}
    \item Undo back to \texttt{"("}. Close allowed ($0<1$): append
          \texttt{)}. $\textsc{Solve}(\texttt{"()"},1,1)$: open allowed
          ($1<2$): append \texttt{(}.
          $\textsc{Solve}(\texttt{"()("},2,1)$: close allowed
          ($1<2$): append \texttt{)}.
          $\textsc{Solve}(\texttt{"()()"},2,2)$: record
          \texttt{"()()"}.
  \end{itemize}
\end{itemize}

Recorded strings: \texttt{"(())"}, \texttt{"()()"} -- exactly the $2$
valid parenthesizations of $n=2$ pairs.

\subsection*{C++ implementation}

\begin{lstlisting}
#include <bits/stdc++.h>
using namespace std;

void solve(string& current, int openCount, int closeCount, int n,
           vector<string>& result) {
    if ((int)current.size() == 2 * n) {
        result.push_back(current);
        return;
    }
    if (openCount < n) {
        current.push_back('(');
        solve(current, openCount + 1, closeCount, n, result);
        current.pop_back();
    }
    if (closeCount < openCount) {
        current.push_back(')');
        solve(current, openCount, closeCount + 1, n, result);
        current.pop_back();
    }
}

vector<string> generateParenthesis(int n) {
    vector<string> result;
    string current;
    solve(current, 0, 0, n, result);
    return result;
}

int main() {
    for (string& s : generateParenthesis(3)) cout << s << " ";
    cout << endl;
    return 0;
}
\end{lstlisting}

\begin{complexitybox}
\textbf{Time Complexity.} The number of valid strings generated is the
$n$-th Catalan number, $C_n = \binom{2n}{n}/(n+1)$, and each costs
$O(n)$ to build, giving an overall time complexity of
\[
  O\!\left(\frac{4^n}{\sqrt{n}}\right)
\]
(the asymptotic growth rate of the Catalan numbers), substantially
smaller than the $O(2^{2n})$ that generating and filtering all
arrangements would cost.
\textbf{Space Complexity.} $O(n)$ for the recursion depth (the string
grows to length $2n$, but the call depth is also $O(n)$ since two
characters can be added back to back in the worst case, or more
precisely the depth equals the string length built so far, $O(n)$ at the
deepest).
\end{complexitybox}

\begin{takeawaybox}
Whenever a backtracking problem's notion of ``validity'' can be checked
\emph{incrementally} (true/false after every partial choice, not only at
the very end), build that incremental check directly into the recursion
as a guard before each choice, rather than generating complete candidates
and validating them afterward. This is dramatically more efficient
whenever invalid prefixes are common, since entire invalid subtrees are
never explored at all.
\end{takeawaybox}
